% The flip-latch, installed: N slots -> N proxies -> ONE shared selector.
% This is the paper's essential diagram; everything in Sec. 4 refers to it.
%
% Drawn in TikZ (not an imported image) so the figure ships as source in the
% arXiv submission and needs no external toolchain to rebuild.
%
% The picture must make exactly one thing obvious: every proxy reads the SAME
% selector word, so the single store that flips it switches the whole set. If a
% reader of the figure can see that, the mechanism is explained.
\begin{figure}[t]
\centering
\begin{tikzpicture}[
  font=\small,
  slot/.style   = {draw, minimum width=15mm, minimum height=6mm, fill=gray!8},
  proxy/.style  = {draw, minimum width=26mm, minimum height=9mm, fill=blue!6},
  target/.style = {draw, rounded corners=2pt, minimum width=11mm, minimum height=5.5mm},
  old/.style    = {target, fill=gray!15},
  new/.style    = {target, fill=green!12},
  sel/.style    = {draw, very thick, minimum width=17mm, minimum height=8mm, fill=orange!18},
  arr/.style    = {-latex, thin},
  dep/.style    = {-latex, thin, densely dotted, gray!70!black},
]

% ---- slots (the words the reader loads) -------------------------------------
\node[slot] (s1) at (0, 1.6)  {\code{s$_1$}};
\node[slot] (s2) at (0, 0)    {\code{s$_2$}};
\node[slot] (s3) at (0, -1.6) {\code{s$_3$}};
\node[above=1.5mm of s1, align=center, font=\footnotesize\itshape] {slots\\(tagged)};

% ---- proxies ---------------------------------------------------------------
\node[proxy] (p1) at (3.6, 1.6)  {$\{o_1,\, n_1\}$};
\node[proxy] (p2) at (3.6, 0)    {$\{o_2,\, n_2\}$};
\node[proxy] (p3) at (3.6, -1.6) {$\{o_3,\, n_3\}$};
\node[above=1.5mm of p1, align=center, font=\footnotesize\itshape] {proxies};

\draw[arr] (s1) -- (p1);
\draw[arr] (s2) -- (p2);
\draw[arr] (s3) -- (p3);

% ---- the shared selector: the whole point ----------------------------------
\node[sel] (sel) at (7.9, 0) {\code{sel}: $0 \to 1$};
\node[above=1.5mm of sel, align=center, font=\footnotesize\itshape] {flip group\\(one word)};

\draw[dep] (p1.east) -- (sel.north west);
\draw[dep] (p2.east) -- (sel.west);
\draw[dep] (p3.east) -- (sel.south west);

% ---- resolution: the WHOLE old set, or the WHOLE new set ------------------
% Not one row expanded and the other two implied -- that misses the point of
% the figure, which is that the selector moves the entire set together. Two
% arrows, two groups; the reader sees three slots resolving as one.
\node[old] (o1) at (11.3,  2.45) {$o_1$};
\node[old] (o2) at (11.3,  1.75) {$o_2$};
\node[old] (o3) at (11.3,  1.05) {$o_3$};
\node[draw, densely dashed, gray!60, inner sep=2.2mm, fit=(o1)(o2)(o3)] (oldset) {};

\node[new] (n1) at (11.3, -1.05) {$n_1$};
\node[new] (n2) at (11.3, -1.75) {$n_2$};
\node[new] (n3) at (11.3, -2.45) {$n_3$};
\node[draw, densely dashed, green!45!black, inner sep=2.2mm, fit=(n1)(n2)(n3)] (newset) {};

% The selector value labels the GROUP, not the arrow. On the path a label is
% crossed by its own line; beside the path it collides with the "flip group"
% caption. On the group it is unambiguous (the arrow points at the group the
% condition selects) and it cannot collide with anything.
\draw[arr, gray!70!black] (sel.east) -- (oldset.west);
\draw[arr, thick]         (sel.east) -- (newset.west);
\node[font=\scriptsize, anchor=south] at (oldset.north) {\code{sel}=0};
\node[font=\scriptsize, anchor=south] at (newset.north) {\code{sel}=1};

\node[align=center, font=\footnotesize\itshape, anchor=north] at (11.3, -3.05)
  {one word decides\\the whole set};

% ---- the flip --------------------------------------------------------------
\draw[-latex, very thick, orange!70!black]
  (sel.south) ++(0,-0.15) -- ++(0,-0.75)
  node[below, align=center, font=\footnotesize\bfseries, text=black]
  {the commit:\\one release store};

\end{tikzpicture}
\caption{The flip-latch, installed, with three slots standing for $N$. Each slot
holds a \emph{tagged} pointer to a proxy rather than to its target; each proxy
carries $\{\mathit{old},\mathit{new}\}$ and reads one shared selector word.
While \code{sel} is $0$ every proxy in the group resolves to its old target, so
installation is invisible to readers. The commit is a single release store,
$0 \to 1$, after which every proxy resolves to its new target. The two dashed
groups are what one store chooses between: there is no instant at which the set
is part-old and part-new, and no per-slot work at commit time.}
\label{fig:fliplatch}
\end{figure}
